\suppnote{Three-city anisotropy results}{note:anisotropy-results}

Median response anisotropy is 0.635, 0.651 and 0.705 in Porto, San Francisco and Chengdu, against structural-tensor medians 0.144, 0.057 and 0.280. The response exceeds structural anisotropy at 79.3\% of Porto cells, 71--72\% of San Francisco cells and 72.4\% of Chengdu cells. These are properties of the fitted response on its identified source span, not of the instantaneous flow distribution. Figure~\ref{fig:supp-4} shows the full distributions behind these medians. Its zero-response atoms are distinct from low but nonzero anisotropy, and its rank-one atoms explain the vertical approach at $A=1$.

Strong local anisotropy does not imply one city-wide preferred direction. Unweighted nematic coherence is 0.087 in Porto, 0.21--0.22 across San Francisco folds and 0.057 in Chengdu. Principal axes vary across space rather than collapsing onto one compass orientation.

Using the joint $\AC\AU$-weighted estimand in Eq.~\ref{eq:supp-axis-agreement}, structural-axis similarity is 0.79, 0.79 and 0.76 in Porto, San Francisco and Chengdu. San Francisco's anisotropy map is reproducible across five folds, with mean pairwise spatial correlation 0.984 (range 0.974--0.991) and anisotropy-weighted nematic axis agreement 0.993. The invariance checks concern scalar anisotropy and covariant axial/tensor transformations as specified in Note~\ref{note:spatial}; they do not imply an axis that remains numerically unchanged under a rigid physical rotation.
